\section{ARSITEKTUR SISTEM}

\subsection{Diagram Arsitektur Sistem}

\begin{figure}[H]
\centering
% Placeholder untuk diagram
\fbox{\parbox{0.8\textwidth}{\centering\vspace{4cm}Diagram Use Case\\(Akan diganti dengan diagram aktual)\vspace{4cm}}}
\caption{Use Case Diagram Savora}
\label{fig:usecase}
\end{figure}

Terdapat tiga aktor utama dalam sistem Savora:
\begin{enumerate}[leftmargin=*]
    \item \textbf{Customer} — Pengguna yang membeli makanan surplus, memberikan rating, dan melacak pesanan.
    \item \textbf{UMKM} — Pengelola produk yang mengelola stok, menentukan harga, dan memproses pesanan.
    \item \textbf{Admin} — Pengelola sistem yang memantau aktivitas, mengelola user, dan memverifikasi donasi.
\end{enumerate}

\subsection{ERD (Entity Relationship Diagram)}

\begin{figure}[H]
\centering
% Placeholder untuk ERD
\fbox{\parbox{0.8\textwidth}{\centering\vspace{4cm}ERD Savora\\(Akan diganti dengan ERD aktual)\vspace{4cm}}}
\caption{Entity Relationship Diagram Savora}
\label{fig:erd}
\end{figure}

\subsubsection{Tabel Utama}

\begin{table}[H]
\centering
\caption{Struktur Database Savora}
\label{tab:database}
\small
\begin{tabularx}{\textwidth}{|l|X|}
\hline
\textbf{Tabel} & \textbf{Deskripsi} \\
\hline
users & Data pengguna (admin, UMKM, customer) \\
\hline
products & Informasi produk makanan \\
\hline
orders & Data transaksi pesanan \\
\hline
order\_items & Detail item dalam pesanan \\
\hline
food\_eligibility & Data kelayakan makanan \\
\hline
reviews & Rating dan ulasan konsumen \\
\hline
transactions & Log pembayaran \\
\hline
donations & Data donasi makanan \\
\hline
notifications & Notifikasi pengguna \\
\hline
\end{tabularx}
\end{table}

\subsection{API Design}

Berikut adalah endpoint API utama yang digunakan dalam Savora:

\begin{table}[H]
\centering
\caption{Endpoint API Utama}
\label{tab:api}
\small
\begin{tabularx}{\textwidth}{|l|l|X|}
\hline
\textbf{Method} & \textbf{Endpoint} & \textbf{Deskripsi} \\
\hline
POST & /api/auth/register & Registrasi pengguna baru \\
\hline
POST & /api/auth/login & Login dan dapatkan token \\
\hline
GET & /api/products & Dapatkan daftar produk \\
\hline
POST & /api/products & Tambah produk baru (UMKM) \\
\hline
PUT & /api/products/:id & Update produk \\
\hline
DELETE & /api/products/:id & Hapus produk \\
\hline
POST & /api/orders & Buat pesanan baru \\
\hline
GET & /api/orders/:id & Dapatkan detail pesanan \\
\hline
PUT & /api/orders/:id/status & Update status pesanan \\
\hline
POST & /api/reviews & Tambah rating \& ulasan \\
\hline
GET & /api/admin/dashboard & Dapatkan data dashboard \\
\hline
\end{tabularx}
\end{table}

\subsection{Alur Data}

Alur data utama dalam Savora:

\begin{enumerate}[leftmargin=*]
    \item \textbf{Alur Penjualan} — UMKM $\rightarrow$ Upload Produk $\rightarrow$ Food Eligibility Assessment $\rightarrow$ Dynamic Pricing $\rightarrow$ Customer Browse $\rightarrow$ Checkout $\rightarrow$ Payment $\rightarrow$ Tracking $\rightarrow$ Delivery $\rightarrow$ Rating.
    
    \item \textbf{Alur Donasi (Manual)} — Food Rescue Score Menurun $\rightarrow$ Notifikasi ke UMKM $\rightarrow$ UMKM Putuskan Donasi $\rightarrow$ Pilih Penerima $\rightarrow$ Koordinasi Pickup.
    
    \item \textbf{Alur Monitoring} — Admin $\rightarrow$ Dashboard $\rightarrow$ Lihat Statistik $\rightarrow$ Identifikasi Issue $\rightarrow$ Tindak Lanjut (suspend, verifikasi, moderasi).
\end{enumerate}
